\documentclass[UTF8,11pt,a4paper]{ctexart}

% ==============================================================================
% 宏包加载与基础配置
% ==============================================================================
\usepackage[top=2.2cm, bottom=2.2cm, left=2.3cm, right=2.3cm]{geometry}
\usepackage{amsmath, amssymb, amsthm, mathtools}
\usepackage{booktabs, tabularx, multirow}
\usepackage{graphicx}
\usepackage{xcolor}
\usepackage{tikz}
\usetikzlibrary{arrows.meta, positioning, shapes.geometric, shapes.misc, calc, trees, backgrounds, fit}
\usepackage{tcolorbox}
\tcbuselibrary{skins, breakable}
\usepackage{enumitem}
\usepackage{fancyhdr}
\usepackage{lastpage}
\usepackage{titlesec}
\usepackage{caption}
\usepackage[hidelinks, colorlinks=true, linkcolor=navyblue, citecolor=navyblue, urlcolor=navyblue]{hyperref}

% ==============================================================================
% 色彩与版面定义
% ==============================================================================
\definecolor{navyblue}{RGB}{27, 54, 93}
\definecolor{slateblue}{RGB}{44, 82, 130}
\definecolor{charcoal}{RGB}{45, 55, 72}
\definecolor{darkgray}{RGB}{74, 85, 104}
\definecolor{lightbg}{RGB}{247, 250, 252}
\definecolor{bordercolor}{RGB}{226, 232, 240}
\definecolor{accentamber}{RGB}{180, 83, 9}

\linespread{1.22}
\setlength{\parskip}{0.32em}
\setlength{\parindent}{2em}

% 标题格式定制
\titleformat{\section}{\color{navyblue}\Large\bfseries}{\thesection}{0.7em}{}
\titleformat{\subsection}{\color{slateblue}\large\bfseries}{\thesubsection}{0.7em}{}
\titleformat{\subsubsection}{\color{charcoal}\normalsize\bfseries}{\thesubsubsection}{0.7em}{}

\titlespacing*{\section}{0pt}{1.0ex plus .2ex minus .2ex}{0.6ex plus .1ex}
\titlespacing*{\subsection}{0pt}{0.8ex plus .2ex minus .1ex}{0.4ex plus .1ex}
\titlespacing*{\subsubsection}{0pt}{0.6ex plus .1ex minus .1ex}{0.3ex plus .1ex}

% 页眉页脚配置
\setlength{\headheight}{14pt}
\pagestyle{fancy}
\fancyhf{}
\fancyhead[L]{\small\color{darkgray}企业大模型使用数据差分隐私分析系统技术方案}
\fancyhead[R]{\small\color{darkgray}Private LLM Analytics}
\fancyfoot[L]{\small\color{darkgray}内部评审方案 \quad\textbar\quad 版本: v1.1-dev}
\fancyfoot[R]{\small\color{darkgray}第 \thepage\ 页\,/\,共 \pageref{LastPage}\ 页}
\renewcommand{\headrulewidth}{0.4pt}
\renewcommand{\footrulewidth}{0.4pt}

% 自定义高亮框
\newtcolorbox{techbox}[2][]{
  enhanced,
  breakable,
  colback=lightbg,
  colframe=navyblue,
  coltitle=white,
  fonttitle=\bfseries\small,
  title={#2},
  boxrule=0.7pt,
  arc=2pt,
  left=8pt,
  right=8pt,
  top=4pt,
  bottom=4pt,
  #1
}

% ==============================================================================
% 封面与元数据配置
% ==============================================================================
\title{
  \vspace{-1.2cm}
  {\LARGE\bfseries\color{navyblue} 企业大模型使用数据差分隐私分析系统技术方案}\\[0.25cm]
  {\large\bfseries\color{slateblue} Private LLM Analytics: Technical Architecture and Privacy-Preserving Analytics Design}\\[0.15cm]
  {\normalsize\color{darkgray} 面向企业内部治理、采购决策与跨组织合规发布的高可用差分隐私解决方案}
}

\author{
  \textbf{系统架构与数据安全联合项目组} \quad\textbar\quad \textit{Private LLM Analytics Engineering Group}
}

\date{
  \small \textbf{版本号：} v1.1-dev \quad\textbar\quad \textbf{发布日期：} 2026 年 10 月 \quad\textbar\quad \textbf{文档状态：} 正式评审方案
  \vspace{-0.6cm}
}

% ==============================================================================
% 正文开始
% ==============================================================================
\begin{document}

\maketitle
\thispagestyle{empty}

\vspace{0.1cm}
\begin{abstract}
随着多模态大语言模型（OpenAI、Kimi、DeepSeek 等）在企业研发、办公与客服场景的深度普及，科学评估员工使用偏好、业务任务结构、Token 成本及接口成功率成为企业数字化治理的关键抓手。然而，传统基于明细日志或简单脱敏均值的统计方式，极易被攻击者结合背景知识通过高维特征反推出特定员工的隐私行为轨迹；而长尾极端调用量用户的存在，更使得统计大盘面临严重的失真与推断风险。

本文系统阐述了 \textbf{\mbox{Private LLM Analytics}} 系统的完整技术方案。系统确立了以单个员工全部行为为边界的“用户级”（\mbox{User-Level DP}）差分隐私数学基石，构建了涵盖调用总量（Total Requests）、Token 消耗（Total Tokens）、员工渗透率（User Prevalence）、接口成功率（Success Rate）四项核心统计能力及扩展的高频任务（Top-K Tasks）分析能力；针对成功率等敏感比率设计了单次联合高斯释放与安全阈值抑制机制；在系统全局治理层面，采用 Rényi 差分隐私（RDP）实现持久化事务一致记账，研发基于规范语义哈希的结果标识缓存重放技术以降低重复查询带来的噪声平均风险；针对长周期监控需求，引入基于时间二叉树（Binary Time Tree）的持续观测机制，使累积噪声从线性降至对数级。在 50,000 名员工规模的基准评测中，系统将宏观相对误差控制在 $1\% \sim 2\%$ 左右，在提供企业级宏观分析能力的同时，降低统计结果泄露个体行为信息的风险。
\end{abstract}

\vspace{0.2cm}
\begin{center}
\small\textbf{关键词：} 差分隐私；大语言模型治理；用户级贡献截断；解析高斯机制；Rényi 差分隐私；持续观测；时间二叉树
\end{center}

\newpage
{
  \setlength{\parskip}{0pt}
  \linespread{1.05}\selectfont
  \renewcommand{\contentsname}{\vspace*{-0.8cm}\color{navyblue}\Large\bfseries 目\quad 录}
  \setcounter{tocdepth}{2}
  \tableofcontents
  \thispagestyle{empty}
}
\newpage

% 正文 10 个章节
\section{建设背景与目标}
\label{sec:background_goals}

\subsection{建设背景：企业大模型应用普及与数据治理需求}
近年来，生成式大语言模型（LLM）逐步融入企业日常经营。企业员工在软件研发、智能客服、文档翻译、合规审查及文案创作等场景中，广泛接入并使用多模态与代码大模型（例如 OpenAI、Kimi、DeepSeek 等）。

为了科学评估各业务部门的采纳情况、合理规划接口采购预算、优化模型负载调度，并对外部模型供应商的服务质量进行客观评估，企业数据治理团队需要定期统计并分析大模型的使用数据。典型的分析需求包括：
\begin{itemize}[leftmargin=2em]
  \item \textbf{部门渗透与使用广度：} 哪些业务部门已经开始使用大模型？日常活跃员工占比如何？
  \item \textbf{交互总量与成本消耗：} 各模型供应商的调用频次、Token 消耗总量及长周期增长趋势如何？
  \item \textbf{任务类型分布：} 员工主要利用模型解决哪些业务任务（如代码编写、翻译、文档总结）？
  \item \textbf{服务可用性与成功率：} 各供应商在复杂任务下的接口调用成功率与异常率如何？
\end{itemize}

\subsection{核心风险：传统统计方式的隐私与数据安全隐患}
在传统数据分析模式下，企业通常直接基于网关明细日志或聚合数据库生成报表。然而，针对大模型交互数据，直接发布明细数据或传统统计均值存在以下隐私与数据安全风险：
\begin{enumerate}[label=\textbf{\arabic*.}, leftmargin=2em]
  \item \textbf{高维行为特征可能导致个体被重新识别：} 大模型调用日志不仅包含员工标识，还记录了调用时间戳、任务类别、Token 长度、使用模型等丰富特征。即使隐去姓名工号进行脱敏，若结合外部公开信息（如考勤打卡时间、特定业务上线节点），依然存在通过高维特征碰撞反推特定员工行为轨迹的可能。
  \item \textbf{极值用户与长尾分布易扭曲大盘统计：} 企业内员工的使用习惯具有显著差异。少数编写自动化脚本的工程师单日可能产生数千甚至数万次调用，而大量普通员工单日仅有数次交互。如果不对单用户影响进行限制，极少数高频用户的行为会主导统计结果，且通过对比特定时段的数据变化，可能推断该员工是否参与了敏感项目。
  \item \textbf{外部数据共享中的业务机密保护：} 当企业向模型供应商反馈质量与成功率指标时，如果直接发布未经保护的细分数据，供应商可能间接推断出企业的核心业务规模与技术选型倾向。
\end{enumerate}

\subsection{建设目标与系统核心定位}
为兼顾“宏观数据可用”与“员工个体隐私”，\textbf{\mbox{Private LLM Analytics}} 系统被设计为企业大模型数据分析与合规发布的统一中间件。系统确立了以下核心建设目标：
\begin{itemize}[leftmargin=2em]
  \item \textbf{数学可证明的隐私保护：} 避免依赖脆弱的静态脱敏规则，引入差分隐私（Differential Privacy）作为底层机制。系统采用以单员工历史行为为边界的用户级保护（User-Level DP），在数学边界内有效限制单个员工记录对输出结果的可辨识影响。
  \item \textbf{高可用业务统计支持：} 在保障隐私预算的前提下，针对企业核心关注的调用量、Token、渗透率、成功率四项核心统计能力及扩展的高频任务等进行算法优化，使宏观统计误差保持在工程可接受的范围内，满足容量规划、财务对账与选型决策需求。
  \item \textbf{全生命周期预算闭环管理：} 建立企业级持久化隐私账本，跟踪控制跨时间、跨指标的累计隐私损耗；提供查询结果复用机制，降低重复查询带来的噪声平均风险。
  \item \textbf{多角色隔离与安全合规发布：} 建立供应商视角隔离与安全阈值抑制机制，支持安全导出供外部生态交互的合规分析报表。
\end{itemize}

\section{系统总体方案}
\label{sec:overall_solution}

\subsection{端到端业务处理全流程}
\textbf{\mbox{Private LLM Analytics}} 系统的核心设计思想是构建一条单向流转、前向封闭、不可逆推的隐私计算管道。整个系统从原始模型交互日志接入，到最终对外输出高可用合规分析结果，遵循如图~\ref{fig:overall_pipeline} 所示的标准业务流。

\begin{figure}[htbp]
\centering
\begin{tikzpicture}[
  scale=0.92, every node/.style={scale=0.92},
  stagebox/.style={draw=navyblue, fill=lightbg, rounded corners=3pt, thick, align=center, inner sep=4pt, text width=3.3cm, minimum height=1.1cm},
  widebox/.style={draw=navyblue, fill=lightbg, rounded corners=3pt, thick, align=center, inner sep=4pt, text width=3.8cm, minimum height=1.1cm},
  arrow/.style={-{Latex[length=2.5mm]}, thick, color=navyblue}
]

  % Row 1: Stages 1 to 4
  \node[stagebox] (S1) at (0, 2.0) {\textbf{1. 原始数据接入}\\[2pt]\footnotesize 日志模式与指纹校验};
  \node[stagebox] (S2) at (4.2, 2.0) {\textbf{2. 用户状态聚合}\\[2pt]\footnotesize 离散记录汇聚为向量};
  \node[stagebox] (S3) at (8.4, 2.0) {\textbf{3. 用户贡献限制}\\[2pt]\footnotesize 单用户范数投影截断};
  \node[stagebox] (S4) at (12.6, 2.0) {\textbf{4. 差分隐私加噪}\\[2pt]\footnotesize 解析高斯扰动注入};

  % Row 2: Stages 5 to 7
  \node[widebox] (S5) at (2.2, 0) {\textbf{5. 隐私预算检查}\\[2pt]\footnotesize 持久化账本事务扣减};
  \node[widebox] (S6) at (6.8, 0) {\textbf{6. 结果合规处理}\\[2pt]\footnotesize 安全阈值抑制与归一};
  \node[widebox] (S7) at (11.4, 0) {\textbf{7. 多角色安全发布}\\[2pt]\footnotesize 看板展示与脱敏导出};

  % Connecting arrows - Row 1
  \draw[arrow] (S1) -- (S2);
  \draw[arrow] (S2) -- (S3);
  \draw[arrow] (S3) -- (S4);

  % Inter-row connector
  \draw[arrow] (S4.south) -- ++(0,-0.45) -| (S5.north);

  % Connecting arrows - Row 2
  \draw[arrow] (S5) -- (S6);
  \draw[arrow] (S6) -- (S7);

\end{tikzpicture}
\caption{Private LLM Analytics 系统端到端处理流水线}
\label{fig:overall_pipeline}
\end{figure}

\subsection{关键处理环节定位与协同逻辑}
系统的七大业务处理阶段各司其职、严密协同，确保了数据隐私与统计可用性的平衡：

\begin{enumerate}[label=\textbf{\arabic*.}, leftmargin=2em]
  \item \textbf{原始数据接入与模式校验：} 接收来自企业统一网关或遥测系统的日志（CSV/Parquet 格式）。系统首先验证数据模式合规性，并流式计算整份数据集的 SHA-256 哈希生成版本指纹，确保分析结果与物理数据快照建立确定性绑定。
  \item \textbf{用户级状态矩阵聚合：} 系统以单个员工（唯一匿名用户标识）为主键，将该员工在观测周期内的全部调用行为投影为状态矩阵（各供应商、各任务上的调用次数与成功次数），为后续实施用户级差分隐私提供基础。
  \item \textbf{用户级贡献限制：} 针对长尾高频用户，系统通过数学投影将其对统计结果的总影响截断在预设的上限安全范数 $C$ 以内，有效抑制极少数高频调用对大盘的支配性影响，使敏感度保持在常数水平。
  \item \textbf{专属差分加噪保护：} 系统根据各统计指标的理论敏感度与预设隐私参数（$\varepsilon, \delta$），通过解析高斯机制数值求解符合条件的最小噪声方差并向全局聚合向量注入高斯噪声，在淡化个体特征的同时保留宏观业务特征。
  \item \textbf{隐私预算控制与审计记账：} 在正式发布前，系统将本次查询转录为 Rényi 差分隐私（RDP）开销，并在持久化 SQLite 账本中执行原子比对与扣减。若该接收方的剩余预算不足，系统实行硬性阻断，防止预算被超额透支。
  \item \textbf{业务后处理与安全阈值抑制：} 对加噪后的数值执行合理的业务规则约束（例如：调用量非负、比例和为 100\%、成功数不超过总调用数）。对于样本量较小的统计单元（加噪后的有效分母低于配置阈值 $\tau$），实施阈值抑制并返回空值，防止因小分母放大引起误导。
  \item \textbf{多角色安全发布：} 面向内部 AI 治理团队提供包含置信区间的可视化大屏；面向外部模型供应商提供权限隔离的结构化报表，实现安全可控的跨组织数据流通。
\end{enumerate}

\section{系统架构设计}
\label{sec:system_architecture}

\subsection{五层总体架构体系}
为确保系统在企业级环境下的稳定性、可扩展性与数学严密性，\textbf{\mbox{Private LLM Analytics}} 采用了分工明确、松耦合的五层架构体系。架构层级如图~\ref{fig:system_architecture} 所示。

\begin{figure}[htbp]
\centering
\begin{tikzpicture}[
  scale=0.88, every node/.style={scale=0.88},
  layerbox/.style={draw=navyblue, fill=lightbg, rounded corners=4pt, thick, inner sep=6pt, text width=14.8cm},
  subbox/.style={draw=slateblue, fill=white, rounded corners=2pt, font=\small, align=center, inner sep=3pt, minimum height=0.75cm},
  arrow/.style={-{Latex[length=2.5mm]}, thick, color=navyblue}
]

  % Layer 1: Ingestion
  \node[layerbox] (L1) at (0, 0) {
    \textbf{\color{navyblue} 1. 数据接入层 (Data Ingestion Layer)}\\[3pt]
    \centering
    \begin{tikzpicture}
      \node[subbox, text width=4.5cm] at (0,0) {日志模式与字段类型校验};
      \node[subbox, text width=4.5cm] at (4.9,0) {全数据流式 SHA-256 指纹};
      \node[subbox, text width=4.5cm] at (9.8,0) {动态多模型与任务分类对齐};
    \end{tikzpicture}
  };

  % Layer 2: Protection
  \node[layerbox, below=0.35cm of L1] (L2) {
    \textbf{\color{navyblue} 2. 隐私保护层 (Privacy Preservation Layer)}\\[3pt]
    \centering
    \begin{tikzpicture}
      \node[subbox, text width=4.5cm] at (0,0) {用户状态矩阵聚合};
      \node[subbox, text width=4.5cm] at (4.9,0) {单用户贡献范数截断};
      \node[subbox, text width=4.5cm] at (9.8,0) {解析高斯加噪标定};
    \end{tikzpicture}
  };

  % Layer 3: Analytics
  \node[layerbox, below=0.35cm of L2] (L3) {
    \textbf{\color{navyblue} 3. 统计分析层 (Analytics Profile Layer)}\\[3pt]
    \centering
    \begin{tikzpicture}
      \node[subbox, text width=2.5cm] at (0,0) {调用总量统计};
      \node[subbox, text width=2.5cm] at (2.85,0) {Token 消耗统计};
      \node[subbox, text width=2.5cm] at (5.7,0) {员工渗透率分析};
      \node[subbox, text width=2.5cm] at (8.55,0) {调用成功率评估};
      \node[subbox, text width=3.2cm] at (11.75,0) {热门任务分析 (扩展)};
    \end{tikzpicture}
  };

  % Layer 4: Release
  \node[layerbox, below=0.35cm of L3] (L4) {
    \textbf{\color{navyblue} 4. 结果发布层 (Result Release Layer)}\\[3pt]
    \centering
    \begin{tikzpicture}
      \node[subbox, text width=4.5cm] at (0,0) {安全阈值抑制过滤};
      \node[subbox, text width=4.5cm] at (4.9,0) {规范重放缓存引擎};
      \node[subbox, text width=4.5cm] at (9.8,0) {合规脱敏结构化导出};
    \end{tikzpicture}
  };

  % Layer 5: Governance & Audit
  \node[layerbox, below=0.35cm of L4] (L5) {
    \textbf{\color{navyblue} 5. 权限审计层 (Security \& Audit Layer)}\\[3pt]
    \centering
    \begin{tikzpicture}
      \node[subbox, text width=4.5cm] at (0,0) {多角色数据视角隔离};
      \node[subbox, text width=4.5cm] at (4.9,0) {持久化隐私预算账本};
      \node[subbox, text width=4.5cm] at (9.8,0) {访问鉴权与安全审计日志};
    \end{tikzpicture}
  };

  % Connections
  \draw[arrow] (L1) -- (L2);
  \draw[arrow] (L2) -- (L3);
  \draw[arrow] (L3) -- (L4);
  \draw[arrow] (L4) -- (L5);

\end{tikzpicture}
\caption{Private LLM Analytics 系统五层分层架构图}
\label{fig:system_architecture}
\end{figure}

\subsection{各层级核心职责与交互机制}
五层架构自底向上交互明确，职责边界清晰：

\begin{enumerate}[label=\textbf{\arabic*.}, leftmargin=2em]
  \item \textbf{数据接入层：}
  负责对接企业内部统一 API 网关或数据湖。系统具备高效的向量化解析引擎，支持百万级规模日志的载入。该层对日志列模式（用户标识、时间戳、模型名称、任务分类、成功标记、Token 计数）进行校验，自动完成维度对齐（支持任意 $M$ 个供应商与 $K$ 个任务类别的动态分类学映射），并流式生成数据集 SHA-256 唯一指纹。
  
  \item \textbf{隐私保护层：}
  系统的核心数学计算引擎。负责将原始明细映射为用户状态矩阵并执行贡献限制（$L_2$ 范数截断），计算各指标的理论敏感度；根据配置的隐私目标（$\varepsilon, \delta$），通过解析高斯机制数值标定所需噪声；通过系统内部随机数模块注入扰动，并提供基于时间二叉树的分层复用能力。在生产环境中，系统禁止接收外部调用方传入的随机种子，有效防范利用已知种子预测噪声的安全风险。
  
  \item \textbf{统计分析层：}
  承载企业核心统计业务逻辑。为每一类指标定义专属的数学处理算子：调用总量与 Token 总量使用全域求和与线性投影；员工渗透率基于去重指示向量与公开员工基数进行比率折算；成功率采用成功数与总请求数的联合加噪；高频任务作为扩展能力，采用词表加噪与安全排序。
  
  \item \textbf{结果发布层：}
  负责加噪结果的业务合规转换与对外输出。针对加噪导致的轻微负数进行截断，对概率分布执行归一化修复；实施安全阈值过滤（当加噪后的有效分母低于配置阈值 $\tau$ 时自动抑制输出为空值，避免除以极小值产生极端偏差）；维护规范重放缓存（ReleaseKey），对历史相同语义的查询直接返回此前发布的受保护结果，降低重复查询造成的无谓预算流失。
  
  \item \textbf{权限审计层：}
  保障数据安全发布的治理支撑。建立以角色为核心的作用域隔离机制：企业内部 AI 治理团队拥有跨模型全景视角；外部模型供应商仅能查询自身模型维度的脱敏指标；维护事务一致的 SQLite 隐私账本，记录每一次发布的时间戳、目标接收方、消耗的 RDP 预算及数据集版本，提供可追溯的安全合规审计。
\end{enumerate}

\section{业务统计分析能力}
\label{sec:analytics_capabilities}

系统目前提供四项核心统计能力，并提供 Top-K 热门任务分析作为扩展能力。针对每个指标，系统均采用“业务释义 $\to$ 单用户贡献限制 $\to$ 差分加噪 $\to$ 最终输出”的针对性设计，在保证数学严谨性的同时，提供直观易懂的业务分析结论。

\subsection{总调用量统计（Total Requests）}
\begin{itemize}[leftmargin=2em]
  \item \textbf{解决的业务问题：} 掌握企业内部对各个大模型（OpenAI、Kimi、DeepSeek 等）及各类任务的宏观调用频次，辅助 IT 基础设施部门评估网关吞吐负载，并在不同供应商之间进行科学的流量调度。
  \item \textbf{业务示例：} 某研发团队单日发起 12,000 次代码生成请求，运营团队发起 3,000 次文案翻译请求。管理层需要了解全公司各部门对不同模型的真实调用总规模。
  \item \textbf{单用户贡献限制策略：} 在无约束状态下，少数编写自动化脚本的工程师单日可能产生大量调用，若直接统计，该极端用户将显著主导统计大盘。系统对每名员工在观测期内的调用量向量 $x_u \in \mathbb{R}^D$ 实施 $L_2$ 范数截断：
  \begin{equation}
    \bar{x}_u = x_u \cdot \min\left(1, \frac{C_{\text{req}}}{\|x_u\|_2}\right)
  \end{equation}
  生产基准配置通常设定 $C_{\text{req}} = 20 \sim 100$。该约束确保任何单名员工对全局总调用量统计的最大欧氏扰动不超过 $C_{\text{req}}$ 次，既保留绝大多数正常员工的调用记录，又有效抑制极端异常流量对大盘的非对称影响。
  \item \textbf{最终输出与呈现：} 输出各供应商在各业务任务下的差分隐私调用总量估计值（非负整数）。在典型基准仿真测试中，宏观调用总量相对误差通常可控制在 $1\% \sim 2\%$ 左右。
\end{itemize}

\subsection{Token 消耗总量统计（Total Tokens）}
\begin{itemize}[leftmargin=2em]
  \item \textbf{解决的业务问题：} 大模型 API 服务通常按照 Token（词元）数量计费。财务与采购团队需要合理掌握各模型、各任务的 Token 消耗总量，以便进行月度账单核算、供应商对账以及跨部门成本分摊。
  \item \textbf{业务示例：} 部分用户发送数万字的长文档分析请求，而多数用户仅交互短句。管理层需要了解整体 Token 账单的宏观结构。
  \item \textbf{单用户贡献限制策略：} 单用户 Token 消耗的离散度远大于调用次数。系统对单用户在所有单元的 Token 消耗向量执行 $L_2$ 截断，基准阈值设为 $C_{\text{tok}} = 100,000$ Token。当某员工总消耗超过 10 万 Token 时，其向量按比例缩放至范数为 10 万，确保全局 $L_2$ 敏感度不超过 $C_{\text{tok}}$。
  \item \textbf{最终输出与呈现：} 输出按模型与任务划分的 Token 消耗总量估计值，支持 API 采购预算的合理规划，统计相对误差在 $1\% \sim 2\%$ 左右。
\end{itemize}

\subsection{员工渗透率统计（User Prevalence）}
\begin{itemize}[leftmargin=2em]
  \item \textbf{解决的业务问题：} 评估大模型技术在企业内部的“普及广度”。管理层不仅关心调用了多少次，更关心到底有多少比例的员工真正使用大模型，是否存在“少数人高频使用、多数人未触达”的状况。
  \item \textbf{业务示例：} 某业务部门拥有 2,000 名员工，如果仅有 100 名员工发起了 10 万次请求，尽管调用量巨大，但员工渗透率仅为 $5\%$。
  \item \textbf{单用户贡献限制与基准对齐：}
  \begin{enumerate}[label=\alph*., leftmargin=1.5em]
    \item \textbf{去重二值化：} 只要员工在该维度至少发起过 1 次调用，该维度即标记为 1，否则为 0。
    \item \textbf{跨维度限制（$L_0$ 截断）：} 限制单个员工最多在 $k=2 \sim 9$ 个不同的分类单元中被计入有效用户（超过则按活跃度保留前 $k$ 个）。单个用户对所有单元活跃计数总和的最大 $L_2$ 影响上限为 $\Delta_2 \le \sqrt{k}$。
    \item \textbf{公开企业总人数对齐（消除分母偏差）：} 传统做法若直接除以“日志中活跃出现的去重用户数”，会导致明显的抽样分母偏差。系统采用企业人事系统公开备案的总员工编制数 $N_{\text{public}}$（例如 2,000 人）作为计算分母：
    \begin{equation}
      \hat{P}_j = \frac{\max(0, \tilde{A}_j)}{N_{\text{public}}}
    \end{equation}
    其中 $\tilde{A}_j$ 为加噪后的各单元去重活跃人数。该设计有效避免了活跃分母波动带来的系统性偏差。
  \end{enumerate}
  \item \textbf{最终输出与呈现：} 输出百分比形式的员工渗透率（如某模型代码任务渗透率为 $46.8\%$），绝对误差在仿真测试中通常低于 1 个百分点。
\end{itemize}

\subsection{任务调用成功率（Success Rate）}
\begin{itemize}[leftmargin=2em]
  \item \textbf{解决的业务问题：} 客观评估各大模型供应商在实际企业网络环境与复杂 prompt 下的接口稳定性与模型可用度，为 SLA 选型评估提供客观参考。
  \item \textbf{核心技术挑战与解决思路：} 直接经验比率在小分母情况下具有较高单用户敏感性和严重效用不稳定性，若强行直接给比率加噪或对极小分母做除法，极易出现数值发散。系统采用\textbf{联合高斯释放与安全阈值抑制机制}：
  \begin{enumerate}[label=\arabic*., leftmargin=1.5em]
    \item \textbf{联合向量贡献限制：} 提取单用户的成功次数向量 $s_u$ 与总请求次数向量 $n_u$。首先基于总请求向量求缩放因子：
    \begin{equation}
      a_u = \min\left(1, \frac{C_{\text{rate}}}{\|n_u\|_2}\right)
    \end{equation}
    同步缩放成功向量与请求向量：$\bar{s}_u = a_u s_u$，$\bar{n}_u = a_u n_u$。将二者拼接为联合向量 $w_u = [\bar{s}_u ; \bar{n}_u]$。由于成功数天然不大于请求数（$\|s_u\|_2 \le \|n_u\|_2$），数学上严格可证联合向量的 $L_2$ 敏感度满足：
    \begin{equation}
      \Delta_2 \le \sqrt{\| \bar{s}_u \|_2^2 + \| \bar{n}_u \|_2^2} \le \sqrt{C_{\text{rate}}^2 + C_{\text{rate}}^2} = \sqrt{2} C_{\text{rate}}
    \end{equation}
    基准配置设定 $C_{\text{rate}} = 20 \sim 50$。
    \item \textbf{单次联合加噪与后处理抑制：} 仅消耗一次差分隐私预算，同时对全量成功数与总请求数注入校准高斯噪声，得到加噪值 $\tilde{S}_j$ 与 $\tilde{N}_j$。后处理流程如下：
    \begin{align}
      \hat{N}_j &= \max(\tilde{N}_j, 0) \\
      \hat{S}_j &= \operatorname{clip}(\tilde{S}_j, 0, \hat{N}_j)
    \end{align}
    系统设置成功率最小发布分母 $\tau = 1.0$（对应配置中的 \texttt{min\_noisy\_denominator} 参数）。若加噪后的有效分母 $\hat{N}_j < \tau$，系统输出 \textbf{SUPPRESSED（空值）}；仅当 $\hat{N}_j \ge \tau$ 时，才发布比率 $\hat{r}_j = \hat{S}_j / \hat{N}_j$。
  \end{enumerate}
  \item \textbf{最终输出与呈现：} 输出各模型与各任务下的接口成功率，绝对偏差在基准测试中通常在 1 个百分点量级左右。
\end{itemize}

\subsection{高频任务分布统计（Top-K Tasks，扩展能力）}
\begin{itemize}[leftmargin=2em]
  \item \textbf{解决的业务问题：} 作为系统的扩展统计能力，算法模块已在 Profile 机制层实现，可根据企业业务部署版本选择启用。该能力用于发现企业员工使用大模型最频繁的前 $K$ 大类业务场景（如代码生成、翻译、文档总结等），辅助企业搭建专属知识库或定制微调模型。
  \item \textbf{单用户贡献限制与加噪：} 单用户在全任务词表中每个任务的贡献经过频次截断，且最多在 $k=5 \sim 10$ 个候选任务中产生影响（$L_2$ 敏感度为 $\sqrt{k}$）。系统向词典中所有任务的调用频次注入高斯噪声，在加噪后的数值空间中选取前 $K$ 个最大值并重排。
  \item \textbf{最终输出与呈现：} 输出排序后的前 $K$ 大热门任务清单及估计热度频次。在典型测试验证中，Top-5 热门任务列表与真实高频任务具有良好的重合度。
\end{itemize}

\section{差分隐私保护机制}
\label{sec:dp_mechanism}

差分隐私（Differential Privacy, 简称 DP）是公认具有良好数学保障的数据隐私保护技术框架。传统的信息安全手段（如数据脱敏、访问控制）主要保护原始明细不被直接窃取，但无法防范从合法的统计分析结果中间接推断特定个人的敏感信息。\textbf{\mbox{Private LLM Analytics}} 通过三步核心机制，构建了从原始日志到最终大屏发布的防护网。

\subsection{保护基石：为什么必须坚持“用户级”差分隐私？}
在技术落地前，必须明确差分隐私的保护颗粒度。业界主要存在两种定义：
\begin{itemize}[leftmargin=2em]
  \item \textbf{记录级保护（Record-Level DP）：} 将单次 API 调用日志视为受保护的基本单元。
  \item \textbf{用户级保护（\mbox{User-Level DP}）：} 将单个员工在企业内的所有历史调用行为视为一个不可分割的整体单元。
\end{itemize}

如果采用传统的记录级保护，当某名员工仅偶尔发起 1 次调用时能够获得保护；但如果某名员工在统计周期内累计发起了 $k=5,000$ 次请求，根据差分隐私的群体隐私推导法则（Group Privacy），其用户级实际隐私保护界会随记录数 $k$ 发生显著恶化，在 $(\varepsilon, \delta)$-DP 框架下失效率 $\delta$ 项也会成倍放大，使得高频活跃员工的隐私保护效果大幅下降。

因此，\textbf{本系统在底层全面采用用户级邻接关系定义}：无论某位员工产生了 10 次还是 10 万次交互，系统在设定的差分隐私参数范围内，将单个用户数据对最终发布结果分布的影响限制在严格受控的边界内，显著降低外部攻击者推断员工个体行为的风险。

\subsection{第一步：为什么要限制单用户贡献（贡献截断）？}
\begin{itemize}[leftmargin=2em]
  \item \textbf{直观原理：} 在企业真实场景中，极个别高频用户（如编写脚本自动化测试的研发人员）可能单日产生大量调用。如果不对单名员工的最大影响力设定上限，为了掩盖这部分调用量的波动，系统就必须添加极大的随机噪声，导致最终统计大盘严重失真而失去业务参考价值。
  \item \textbf{系统做法（用户级 $L_2$ 范数截断）：} 系统对每位员工聚合后的高维向量 $x_u$ 施加几何投影截断：
  \begin{equation}
    \bar{x}_u = x_u \cdot \min\left(1, \frac{C}{\|x_u\|_2}\right)
    \label{eq:l2_clipping}
  \end{equation}
  其中 $C$ 为预设的安全上限截断阈值（例如调用量 $C_{\text{req}}=20 \sim 100$）。
  \item \textbf{核心效果：} 正常员工的调用向量（$\|x_u\|_2 \le C$）不受缩放影响；而极少数异常高频用户的贡献向量被等比例压缩至欧氏长度恰好为 $C$。这一步骤将统计函数的最大单用户变动量（即全局 $L_2$ 敏感度 $\Delta_2$）严格锚定在常数 $C$ 以内，为后续校准加噪提供了有限且可控的数学前提。
\end{itemize}

\subsection{第二步：为什么要加噪声（解析高斯机制）？}
\begin{itemize}[leftmargin=2em]
  \item \textbf{直观原理：} 限制了单用户贡献后，如果直接发布求和结果，攻击者对比“该员工入职前”与“入职后”两个报表的差值，依然能反算出该员工的截断数据。因此，必须向最终的汇总向量中注入恰当大小的随机扰动（加噪），将单用户贡献淹没在统计波动之中。
  \item \textbf{系统做法（高斯扰动注入）：} 系统在全量截断数据求和后，添加均值为 0、协方差为 $\sigma^2 \mathbf{I}$ 的多维正态高斯噪声：
  \begin{equation}
    \tilde{Y} = \sum_{u=1}^N \bar{x}_u + \mathcal{N}\left(0, \sigma^2 \mathbf{I}\right)
    \label{eq:gaussian_mechanism}
  \end{equation}
  \item \textbf{噪声标定技术（Analytic Gaussian Mechanism）：} 噪声标准差 $\sigma$ 不能凭经验主观设定。加多了影响业务可用性，加少了达不到设定的隐私保护强度。系统集成了 Balle \& Wang 提出的\textbf{解析高斯标定算法}，根据设定的敏感度 $\Delta_2$、隐私损失上限 $\varepsilon$ 和失败概率 $\delta$，数值求解满足差分隐私条件的最小标准差 $\sigma$。相比传统经典高斯上界公式，解析高斯机制有效避免了松弛放缩，减少了多余的噪声添加。
\end{itemize}

\subsection{第三步：为什么要控制总体消耗（隐私预算）？}
\begin{itemize}[leftmargin=2em]
  \item \textbf{直观原理：} 单次发布被噪声保护，但如果针对同一批底层数据发起多次统计查询，多次加噪结果相互组合可能会逐步累积信息泄露。差分隐私将这种累积泄露量量化为\textbf{“隐私预算”}，通常记作 $(\varepsilon, \delta)$。
  \item \textbf{通俗理解隐私参数：}
  \begin{itemize}
    \item $\varepsilon$（隐私损失系数）：衡量数据发布前后推断概率的最大变动倍数。$\varepsilon$ 越小（如 $\varepsilon \le 1.0$），隐私保护越强，噪声相对越大；$\varepsilon$ 适中（如 $\varepsilon = 1.0 \sim 2.0$），在工业场景中实现了良好的效用与安全性平衡。
    \item $\delta$（失效率上限）：极小概率事件的松弛上界，代表差分隐私成立的置信边界。系统默认设置 $\delta = 1.0 \times 10^{-6}$（百万分之一），其概率远低于硬件物理随机故障率，在工程上具有充分的可靠性。
  \end{itemize}
  \item \textbf{全生命周期闭环记账：} 系统将企业允许消耗的总隐私预算视为固定资产。每一次统计发布都必须消耗特定的预算配额。当配额达到上限时，系统执行阻断，拒绝新增查询，有效控制多次查询带来的累积泄露风险。
\end{itemize}

\section{隐私预算与重复查询控制}
\label{sec:budget_and_replay}

在差分隐私工程落地中，仅有优秀的加噪算法是不够的。如果系统缺乏对历史发布的全局管控，恶意分析人员可以通过发起多次重复查询，将多次加噪的结果进行简单平均，从而逐步抵消高斯噪声、还原出原始数据特征（即“噪声消除与重构风险”）。\textbf{\mbox{Private LLM Analytics}} 通过引入 \textbf{RDP 隐私记账机制}与 \textbf{ReleaseKey 规范重放机制}，从工程架构上有效应对此类风险。

\subsection{隐私预算记账机制（RDP Accountant）}
传统的差分隐私累加规则（如基础组合定理 Basic Composition）往往较为保守，若简单按逐次相加计算，多项联合发布或连续发布时会过快耗尽预算，或者迫使系统添加过大噪声。

为此，系统记录每一次新发布带来的隐私消耗，并累计计算总预算；达到配置上限后停止新的发布。底层采用 RDP 记账，在预设的一组 RDP 阶数上累计隐私消耗，并在给定 $\delta$ 下选择对应的最小 $\varepsilon$ 上界，以获得比简单逐次相加更合理的累计预算估计。

在具体工程实现上，系统隐私账本采用轻量级嵌入式 SQLite 存储，支持 ACID 事务特性。在多线程或微服务并发调用时，系统采用“比对预留（Compare-and-Reserve）”原子事务机制，确保并发场景下隐私预算不发生超额扣减或透支。

\subsection{ReleaseKey 规范重放机制：降低重复查询带来的噪声平均风险}
\begin{itemize}[leftmargin=2em]
  \item \textbf{安全考量（噪声平均风险）：} 假设针对某一固定时段的企业调用总量发起 100 次相同的统计请求。若系统每次都重新注入独立的高斯随机噪声，根据统计学大数定律，将这 100 次返回的加噪值求简单平均，其随机噪声的标准差将衰减至原始噪声的 $\frac{1}{\sqrt{100}} = \frac{1}{10}$，导致差分隐私防护水平大幅下降。
  \item \textbf{解决方案（结果标识生成与安全复用）：} 系统根据数据版本、指标、查询参数和接收方等信息生成结果标识（ReleaseKey），相同的已发布查询可直接复用原结果：
  \begin{equation}
    \text{ReleaseKey} = \operatorname{SHA256}\Big(\text{Dataset\_Fingerprint} \parallel \text{Metric} \parallel \text{Params} \parallel \text{Recipient}\Big)
  \end{equation}
  \begin{enumerate}[label=\alph*., leftmargin=1.5em]
    \item \textbf{命中历史查询：} 若该 ReleaseKey 已在持久化账本中存在，系统直接调取历史已发布的受保护结果与置信区间并返回。该过程不重新调用随机加噪引擎，增量隐私预算消耗记为 0。重复查询获得的结果保持一致，有效防范了多次重复查询抵消噪声的风险。
    \item \textbf{查询参数或数据变更：} 若分析人员调整了过滤条件、时间范围或截断范数，或者底层遥测日志文件发生了更新（Dataset SHA-256 改变），系统计算得出全新的 ReleaseKey，正常执行预算扣减与加噪发布流程。
  \end{enumerate}
  \item \textbf{管理收益：} 该机制既防范了恶意统计攻击，又为高并发可视化大屏和日常重复报表刷新节省了宝贵的差分隐私预算资产。
\end{itemize}

\section{持续统计机制}
\label{sec:continual_observation}

在真实企业运营场景中，大模型的使用监控通常是一个长周期、跨天、跨月的\textbf{持续观测过程}。管理层不仅需要查看当天的即时调用量，还需要调阅最近 7 天的滚动周报、最近 30 天的月度趋势，乃至连续按天查看历史增长曲线。

\subsection{长周期观测的隐私预算挑战}
在传统的差分隐私方案中，跨时间统计面临着较为严峻的预算累积挑战：
\begin{itemize}[leftmargin=2em]
  \item \textbf{天级独立发布模式：} 如果每天对当天的调用量发布一次独立加噪结果，根据隐私组合原理，30 天累计消耗的隐私预算将线性增加 30 倍。若要将总预算控制在安全阈值以内，分配给每一天的预算将非常微小，注入的噪声容易对日常小幅度的业务趋势造成干扰。
  \item \textbf{滑动窗口重复查询：} 如果每天重新对“过去 30 天累计调用量”发起一次全量查询，同一批历史数据将被反复加噪释放，不仅加速消耗总预算，且重复查询增加了被统计分析利用的暴露面。
\end{itemize}

\subsection{时间二叉树分层复用机制}
为了兼顾“长期连续监控”与“预算节约”，\textbf{\mbox{Private LLM Analytics}} 实现了基于\textbf{时间二叉树（Binary Time Tree）}的分层统计复用机制。

\begin{figure}[htbp]
\centering
\begin{tikzpicture}[
  scale=0.9, every node/.style={scale=0.9},
  treenode/.style={draw=navyblue, fill=white, rounded corners=3pt, thick, align=center, inner sep=3pt, font=\small},
  hitnode/.style={draw=accentamber, fill=accentamber!20, rounded corners=3pt, very thick, align=center, inner sep=3pt, font=\bfseries\small},
  leafnode/.style={draw=slateblue, fill=lightbg, rounded corners=2pt, align=center, inner sep=3pt, font=\small},
  edge/.style={thick, color=darkgray}
]

  % Level 3 (Root)
  \node[treenode, text width=6.8cm, minimum height=0.75cm] (N_1_8) at (7.5, 3.8) {根节点：区间 $[1, 8]$\\\footnotesize (8天整体汇总)};

  % Level 2 (4-day)
  \node[hitnode, text width=5.0cm, minimum height=0.75cm] (N_1_4) at (3.55, 2.5) {区间 $[1, 4]$\\\footnotesize (命中复用项 1)};
  \node[treenode, text width=5.0cm, minimum height=0.75cm] (N_5_8) at (11.45, 2.5) {区间 $[5, 8]$\\\footnotesize (4天跨度汇总)};

  % Level 1 (2-day)
  \node[treenode, text width=2.6cm, minimum height=0.7cm] (N_1_2) at (1.85, 1.25) {区间 $[1, 2]$};
  \node[treenode, text width=2.6cm, minimum height=0.7cm] (N_3_4) at (5.25, 1.25) {区间 $[3, 4]$};
  \node[hitnode, text width=2.6cm, minimum height=0.7cm] (N_5_6) at (9.75, 1.25) {区间 $[5, 6]$\\\footnotesize (命中复用项 2)};
  \node[treenode, text width=2.6cm, minimum height=0.7cm] (N_7_8) at (13.15, 1.25) {区间 $[7, 8]$};

  % Level 0 (1-day leaves)
  \node[leafnode, text width=1.45cm, minimum height=0.65cm] (L1) at (1.0, 0) {第 1 天};
  \node[leafnode, text width=1.45cm, minimum height=0.65cm] (L2) at (2.7, 0) {第 2 天};
  \node[leafnode, text width=1.45cm, minimum height=0.65cm] (L3) at (4.4, 0) {第 3 天};
  \node[leafnode, text width=1.45cm, minimum height=0.65cm] (L4) at (6.1, 0) {第 4 天};
  \node[leafnode, text width=1.45cm, minimum height=0.65cm] (L5) at (8.9, 0) {第 5 天};
  \node[leafnode, text width=1.45cm, minimum height=0.65cm] (L6) at (10.6, 0) {第 6 天};
  \node[hitnode, text width=1.45cm, minimum height=0.65cm] (L7) at (12.3, 0) {第 7 天\\\footnotesize (项 3)};
  \node[leafnode, text width=1.45cm, minimum height=0.65cm] (L8) at (14.0, 0) {第 8 天};

  % Edges
  \draw[edge] (N_1_8.south) -- (N_1_4.north);
  \draw[edge] (N_1_8.south) -- (N_5_8.north);

  \draw[edge] (N_1_4.south) -- (N_1_2.north);
  \draw[edge] (N_1_4.south) -- (N_3_4.north);
  \draw[edge] (N_5_8.south) -- (N_5_6.north);
  \draw[edge] (N_5_8.south) -- (N_7_8.north);

  \draw[edge] (N_1_2.south) -- (L1.north);
  \draw[edge] (N_1_2.south) -- (L2.north);
  \draw[edge] (N_3_4.south) -- (L3.north);
  \draw[edge] (N_3_4.south) -- (L4.north);
  \draw[edge] (N_5_6.south) -- (L5.north);
  \draw[edge] (N_5_6.south) -- (L6.north);
  \draw[edge] (N_7_8.south) -- (L7.north);
  \draw[edge] (N_7_8.south) -- (L8.north);

\end{tikzpicture}
\caption{持续观测时间二叉树分解示意图（以查询前 7 天区间 $[1, 7]$ 为例）}
\label{fig:time_tree}
\end{figure}

\subsection{时间二叉树的核心运作原理与业务优势}
\begin{enumerate}[label=\textbf{\arabic*.}, leftmargin=2em]
  \item \textbf{树状分层结构：} 如图~\ref{fig:time_tree} 所示，系统将观测时间周期（例如 16 天或 32 天）构建为一棵二叉树。底层叶子节点对应单日数据聚合，父节点对应连续 2 天、4 天、8 天乃至全周期的聚合节点。
  \item \textbf{规范前缀对偶分解：} 任意连续时间区间的统计查询，均可被唯一拆解为该树中最多 $\lceil \log_2 T \rceil$ 个规范互斥节点的代数和。例如，查询“第 1 天至第 7 天”的前缀总调用量，不需要相加 7 个独立单日叶子节点，而是直接分解为 3 个节点之和：
  \begin{equation}
    \text{区间} [1, 7] = \text{节点}[1, 4] + \text{节点}[5, 6] + \text{节点}[7, 7]
  \end{equation}
  \item \textbf{累积噪声随时间呈对数级增长：}
  由于任意时间范围的估计值至多由树高（$\log_2 T$）个加噪节点合成，累积噪声的方差仅随时间呈对数增长（$O(\log T)$），而非线性的 $O(T)$。
  \begin{itemize}
    \item 观测周期达到 32 天时，直接天级累加需要叠加 32 次独立噪声；而时间二叉树仅需合成至多 5 个分层节点，大幅减少了累积噪声对长周期趋势的干扰。
    \item 在连续监控下，跨天重复节点自动复用历史加噪值，有效降低了对企业年度总隐私预算的日常消耗。
  \end{itemize}
  \item \textbf{代际轮换（Epoch Rolling）：} 当一棵二叉树的预设跨度（如 30 天）走完后，系统滚动开启下一个代际树，支持长周期持续安全监控。
\end{enumerate}

\section{权限、安全与结果发布}
\label{sec:security_and_release}

数据要素安全流通的核心在于构建清晰的信任边界。\textbf{\mbox{Private LLM Analytics}} 在架构层面建立了“角色隔离、底层加固、输出阻断”的多维主动安全防护网。

\subsection{多角色视角隔离机制（Recipient Scopes）}
企业在治理大模型数据时，面临内部管理与外部供应商两套截然不同的使用主体。系统在数据源头与发布策略上建立了严格的接收方作用域（Recipient Scope）边界：
\begin{enumerate}[label=\textbf{\arabic*.}, leftmargin=2em]
  \item \textbf{内部治理视角（\texttt{internal}）：} 面向企业内部 AI 治理委员会、安全团队与技术管理层。该角色拥有跨模型、跨部门的全景分析权限，能够同时对比 OpenAI、Kimi、DeepSeek 等所有供应商的调用量、成本占比、渗透率与可用性指标，辅助宏观技术决策。
  \item \textbf{外部供应商隔离视角（\texttt{provider\_a/b/c}）：} 面向外部模型服务商进行服务质量对齐时，系统对数据实行严格的作用域隔离与过滤：
  \begin{itemize}
    \item 针对 OpenAI（\texttt{provider\_a}）发布的报告，仅包含该供应商自身的接口成功率与任务调用量，不包含任何竞品（Kimi、DeepSeek）的统计信息；
    \item 针对 Kimi（\texttt{provider\_b}）与 DeepSeek（\texttt{provider\_c}）同样实施双向隔离。
  \end{itemize}
  \item \textbf{全局外泄控制线（\texttt{GLOBAL\_EXTERNAL}）：} 系统设置了跨外部供应商的合并隐私预算水位线，防止多个外部供应商通过串谋或联合分析拼凑出企业全局业务大盘。
\end{enumerate}

\subsection{生产环境随机数安全与抗操纵设计}
差分隐私的数学有效性依赖于噪声样本的不可预测性。在工程安全审计中，系统确立了严格的防注入规范：
\begin{itemize}[leftmargin=2em]
  \item \textbf{生产模式禁止外部种子注入：} 生产模式禁止调用方指定随机种子（\texttt{seed}）或随机数生成器（\texttt{rng}）。若外部请求试图传入自定义种子，系统触发安全拦截异常（ProductionRNGControlError）并阻断发布流程。
  \item \textbf{系统内部随机数独立生成：} 扰动噪声由系统内部随机数模块独立生成，加噪过程内存隔离，保障扰动采样的统计独立性。
\end{itemize}

\subsection[小分母安全阈值抑制与合规后处理]{小分母安全阈值抑制（$\tau$-Suppression）与合规后处理}
除了高斯加噪外，系统在发布前对加噪结果实施严格的合规后处理校验：
\begin{itemize}[leftmargin=2em]
  \item \textbf{非负截断与总和一致性：} 针对调用量、Token 总量等物理计数指标，系统将加噪产生的偶发负数自动截断为 0；针对比例分布指标，进行和为 $100\%$ 的仿射投影约束。
  \item \textbf{安全阈值抑制（Safety Threshold）：} 当某个统计单元加噪后的有效分母低于配置阈值（系统基准阈值设定为 $\tau = 1.0$，对应配置中的 \texttt{min\_noisy\_denominator} 参数；加噪后的有效分母低于该值时不发布成功率）时，系统将该项成功率输出置为 \textbf{SUPPRESSED（空值）}。该规则既消除了小样本高方差导致的数据失真，又降低了针对孤立个案的推断风险。
\end{itemize}

\section{系统部署与实施方案}
\label{sec:deployment_implementation}

\subsection{部署拓扑与运行环境配置}
\textbf{\mbox{Private LLM Analytics}} 遵循轻量化、高兼容性与无状态计算服务的设计规范，可无缝嵌入企业现有的私有云容器平台（如 \mbox{Kubernetes}、Docker）或内部安全网络隔离区。

\begin{figure}[htbp]
\centering
\begin{tikzpicture}[
  scale=0.88, every node/.style={scale=0.88},
  box/.style={draw=navyblue, fill=lightbg, rounded corners=3pt, thick, align=center, inner sep=6pt, font=\small},
  dbbox/.style={draw=slateblue, fill=white, rounded corners=2pt, thick, align=center, inner sep=5pt, font=\small},
  arrow/.style={-{Latex[length=2.5mm]}, thick, color=navyblue}
]

  % Source zone
  \node[box, text width=3.4cm] (GW) at (0, 0) {企业大模型网关 / 遥测系统\\(Kong / One-API 等)};
  
  % Analytics core zone
  \node[box, text width=5.6cm, draw=slateblue, fill=white, very thick] (APP) at (6.6, 0) {
    \textbf{\color{navyblue} \mbox{Private LLM Analytics} 核心服务}\\[3pt]
    \footnotesize
    $\bullet$ 矢量化解析与用户矩阵聚合\\
    $\bullet$ 专属差分隐私加噪计算引擎\\
    $\bullet$ ReleaseKey 缓存与阈值过滤
  };

  % Ledger
  \node[dbbox, below=1.2cm of APP, text width=4.5cm] (DB) {
    \textbf{持久化隐私账本}\\
    \footnotesize (SQLite / 事务 ACID 审计)
  };

  % Consumers
  \node[box, text width=3.8cm] (INT_DASH) at (13.4, 1.2) {内部治理可视化看板\\[1pt]{\footnotesize (全景业务洞察与审计)}};
  \node[box, text width=3.8cm] (EXT_EXP) at (13.4, -1.2) {外部合规脱敏导出\\[1pt]{\footnotesize (OpenAI / Kimi / DeepSeek)}};

  % Links
  \draw[arrow] (GW) -- node[above=2pt, font=\footnotesize] {遥测日志} (APP);
  \draw[arrow] (APP.east) -- node[above, font=\footnotesize] {全景大屏} (INT_DASH.west);
  \draw[arrow] (APP.east) -- node[below, font=\footnotesize] {合规导出} (EXT_EXP.west);
  \draw[arrow, <->] (APP) -- node[right=3pt, font=\footnotesize] {预留与记账} (DB);

\end{tikzpicture}
\caption{企业大模型差分隐私分析系统部署拓扑与数据流向图}
\label{fig:deployment}
\end{figure}

\begin{itemize}[leftmargin=2em]
  \item \textbf{算力与存储开销低：} 得益于高度优化的 NumPy 矢量化稀疏状态矩阵算子，系统无需昂贵的 GPU 硬件。单台标准通用服务器（8 核 CPU、16 GB 内存）即可在数秒内完成数万名员工、百万级日志的全流程聚合与加噪发布。
  \item \textbf{环境依赖：} 运行环境支持主流 Linux 发行版（如 Ubuntu 20.04+、RHEL 8+）及现代 Windows Server 环境，运行时依赖标准 Python 3.10+，无特殊专有库绑定。
  \item \textbf{状态持久化：} 隐私状态与账本默认采用嵌入式 SQLite（\texttt{privacy\_ledger.db}），支持单文件热备与异地容灾，亦可平滑对接企业级 \mbox{PostgreSQL} 关系型数据库。
\end{itemize}

\subsection{与企业网关及日志系统的集成步骤}
系统与企业现有 IT 资产的对接流程清晰快捷，分为三步实施：
\begin{enumerate}[label=\textbf{第 \arabic* 步：}, leftmargin=4em]
  \item \textbf{确定标准遥测日志输入（Schema 对齐）：} 网关将日常大模型调用日志以 CSV 或 Parquet 格式导出。标准字段仅需包含：匿名员工标识（\texttt{user\_id}）、请求时间（\texttt{timestamp}）、模型供应商（\texttt{provider}）、业务任务标签（\texttt{task\_category}）、是否调用成功（\texttt{is\_success}）以及消耗 Token 数（\texttt{total\_tokens}）。
  \item \textbf{隐私目标与指标策略配置：} 在企业级 YAML 配置文件中，预先审定各指标的差分隐私参数（推荐 $\varepsilon = 1.0 \sim 2.0$，$\delta = 1.0 \times 10^{-6}$）、用户级截断阈值（$C_{\text{req}}, C_{\text{tok}}, C_{\text{rate}}$）以及接收方隔离规则。
  \item \textbf{调度流水线集成与安全导出：} 通过企业现有的调度平台（如 Airflow、DolphinScheduler 或定时任务）触发批处理或准实时计算任务。结果自动通过 RESTful API 提供给内部 BI 看板展示，或生成只包含差分隐私脱敏结果的安全 CSV 文件分发给合作供应商。
\end{enumerate}

\section{实施效果、适用范围与结论}
\label{sec:utility_and_conclusion}

\subsection{统计精度与效用表现}
系统在基于真实业务特征构建的 50,000 名员工基准仿真数据集上进行了效用测试（默认隐私预算 $\varepsilon = 1.0 \sim 2.0$，$\delta = 1.0 \times 10^{-6}$）。典型参考结果如表~\ref{tab:utility_summary} 所示。

\begin{table}[htbp]
\centering
\small
\renewcommand{\arraystretch}{1.08}
\caption{基准规模（50,000 用户）下各核心指标的差分隐私统计参考表现}
\label{tab:utility_summary}
\begin{tabularx}{\textwidth}{p{2.6cm}>{\raggedright\arraybackslash}p{3.8cm}cc>{\raggedright\arraybackslash}X}
\toprule
\textbf{指标} & \textbf{关键参数设定} & \textbf{评估方式} & \textbf{参考误差水平} & \textbf{适用场景说明} \\
\midrule
总调用量 & $C_{\text{req}} = 20 \sim 100$ & 相对误差 MAPE & 约 $1\% \sim 2\%$ & 适用于容量规划与负载调度 \\
Token 消耗 & $C_{\text{tok}} = 100,000$ & 相对误差 MAPE & 约 $1\% \sim 2\%$ & 适用于采购对账与成本分摊 \\
员工渗透率 & $k=2 \sim 9, N_{\text{public}}=2,000$ & 绝对百分点误差 & 低于 1 个百分点 & 适用于普及度观察与采纳分析 \\
调用成功率 & $C_{\text{rate}} = 20 \sim 50, \tau=1.0$ & 绝对百分点误差 & 约 1 个百分点量级 & 适用于供应商稳定性评估 \\
高频任务 (扩展) & $k=5 \sim 10$, 词表加噪 & Top-5 命中率 & 前 5 大任务良好重合 & 适用于热门需求挖掘与定制微调 \\
\bottomrule
\end{tabularx}
\end{table}

差分隐私数学机制的一个重要特性是：\textbf{注入的高斯噪声规模仅取决于单用户截断阈值，而与企业总人数 $N$ 无关}。因此，企业员工规模越大（$N$ 越大），单用户截断与加噪对宏观总和的相对扰动越微小，相对误差呈 $O(1/N)$ 迅速衰减。在万人级企业中，统计大盘精度与未加噪的真实宏观值高度接近。

\subsection{适用场景与边界说明}
为确保方案在企业内的科学应用，明确系统的数据边界如下：
\begin{itemize}[leftmargin=1.8em, itemsep=1pt]
  \item \textbf{推荐适用场景：}
  \begin{itemize}[itemsep=1pt]
    \item 500 人以上规模的中大型企业与技术研发组织；
    \item 跨多模型（OpenAI、Kimi、DeepSeek 等）的宏观负载分析与采购对账；
    \item 向外部云服务商或行业生态反馈服务质量（SLA）等安全共享场景。
  \end{itemize}
  \item \textbf{不适用场景及安全边界：}
  \begin{itemize}[itemsep=1pt]
    \item \textbf{微型团队或稀疏统计单元：} 样本量较小时，差分隐私噪声的相对影响会更加明显；对于成功率等比率指标，当加噪后的有效分母低于配置阈值时，系统会抑制该项结果。
    \item \textbf{个体违规行为调查：} 差分隐私从原理上用于降低个体行为推断风险，排查个案违规应按照正规内控流程调用审计日志。
  \end{itemize}
\end{itemize}

\subsection{方案总结}
\textbf{\mbox{Private LLM Analytics}} 系统采用用户级差分隐私机制，构建了“单用户贡献截断 $\to$ 解析高斯加噪 $\to$ 隐私预算全生命周期记账 $\to$ 结果标识规范重放 $\to$ 分层时间持续监控”的技术架构，在降低员工个体行为推断风险的同时，提供可用的大盘宏观管理洞察，为企业大模型应用的合规治理与数字化决策提供技术支撑。


% 附录
\newpage
\appendix
\section{附录：系统核心配置与安全参数速查}
\label{sec:appendix_parameters}

本附录汇总了 \textbf{\mbox{Private LLM Analytics}} 在标准企业级生产环境中的基准参数配置矩阵，供系统实施与安全审计人员参考。

\begin{table}[htbp]
\centering
\small
\caption{系统核心差分隐私机制与安全参数配置规范表}
\label{tab:parameter_reference}
\begin{tabularx}{\textwidth}{ll>{\raggedright\arraybackslash}X}
\toprule
\textbf{参数类别} & \textbf{配置项名称} & \textbf{生产默认值与业务约束说明} \\
\midrule
\multirow{3}{*}{基础隐私预算} 
  & 目标单次预算 $\varepsilon$ & $1.0 \sim 2.0$（平衡效用与隐私的推荐工业界区间） \\
  & 目标失效率 $\delta$ & $1.0 \times 10^{-6}$（远低于硬件随机故障率） \\
  & 全局生命周期上限 $\varepsilon_{\text{cap}}$ & $4.0$（达到上限自动触发记账熔断保护） \\
\midrule
\multirow{4}{*}{贡献截断阈值} 
  & 调用量截断 $C_{\text{req}}$ & $20.0 \sim 100.0$ 次 / 周期（规范单用户调用量欧氏范数界） \\
  & Token 截断 $C_{\text{tok}}$ & $100,000$ Token / 周期（规范单用户 Token 欧氏范数界） \\
  & 渗透率跨度限制 $k$ & $2 \sim 9$（单员工最多计入 $k$ 个分类，$\Delta_2 \le \sqrt{k}$） \\
  & 成功率联合限制 $C_{\text{rate}}$ & $20.0 \sim 50.0$（联合向量 $w_u = [\bar{s}_u; \bar{n}_u]$，$\Delta_2 \le \sqrt{2}C_{\text{rate}}$） \\
\midrule
\multirow{2}{*}{发布安全规则} 
  & 成功率最小发布分母 $\tau$ & $1.0$（对应 \texttt{min\_noisy\_denominator}；加噪后的有效分母低于该值时不发布成功率） \\
  & 公开基准人数 $N_{\text{public}}$ & $2,000$（基于人事花名册公开总编制，消除分母偏差） \\
\midrule
\multirow{3}{*}{持续观测机制} 
  & 时间二叉树代际跨度 & $16$ 或 $32$ 天（任意区间至多分解为 $\lceil \log_2 T \rceil$ 个规范节点） \\
  & 节点前缀复用 & 自动复用历史加噪节点，累积噪声仅随时间呈对数级增长 \\
  & 代际平滑滚动 & 到期自动生成新代际树，历史节点归档为只读快照 \\
\midrule
\multirow{3}{*}{权限作用域} 
  & 内部治理（\texttt{internal}） & 全指标、全模型可见（OpenAI、Kimi、DeepSeek 统一透视） \\
  & 外部供应商隔离 & \texttt{provider\_a/b/c} 相互隔离，只能查看自身脱敏数据 \\
  & 外部合并上限 & \texttt{GLOBAL\_EXTERNAL} 追踪跨外部供应商的合并外泄 \\
\bottomrule
\end{tabularx}
\end{table}

\begin{techbox}{合规发布核对清单（Compliance Checklist）}
\begin{enumerate}[leftmargin=1.5em, itemsep=2pt]
  \item \textbf{无明细暴露：} 系统向外部或大屏发布的任何数据均经过用户级差分加噪，不包含任何单条明细或未加噪真实计数。
  \item \textbf{重放无增量消耗：} 相同 ReleaseKey 查询直接复用已有受保护结果，增量预算消耗记为 0。
  \item \textbf{生产随机数防注入：} 生产环境下禁止调用方指定外部随机种子或生成器，扰动噪声由系统内部模块独立生成。
  \item \textbf{小分母抑制生效：} 有效分母低于 $\tau=1.0$ 的统计单元均被自动遮蔽为 SUPPRESSED（空值），防范孤立个案推断。
\end{enumerate}
\end{techbox}


\end{document}
